This thesis asked whether early diffusion structure carries useful information beyond early volume and speed, and whether that information differs across two tasks: veracity screening and eventual reach prediction. With the main empirical answers now established in the Results section, the discussion focuses on how those answers should be interpreted, how they connect back to the literature, and where the main boundaries of the evidence lie.

\subsection{What the Results Show}
The clearest pattern in the revised results is still task asymmetry. On FibVID, veracity performance is above chance but modest throughout the main $k$-window analysis. In that sense, early diffusion signals do contain some information about misinformation risk, but they do not support a strong classification claim.

For reach, the pattern is much stronger. Under the harmonized FibVID specification, prediction performance is high even at moderate windows and remains strong in later ones. This suggests that early diffusion structure is much more informative for \emph{how far} a cascade will spread than for \emph{whether} it is false.

The Twitter15/16 benchmark supports the same broad reading. In the original rumor-tree setting, veracity reaches about AUC $=0.657$ at $k=20$, while reach reaches about $R^2=0.719$ at $k=180$. Relative to that benchmark, FibVID weakens the veracity side but strengthens the reach side. The comparison therefore does not overturn the earlier asymmetry; it makes it clearer.

This point also matters for how the results speak back to the literature. Work such as \citet{vosoughi_spread_2018} shows that true and false information can diffuse differently when cascades are observed in full. The present results suggest a narrower conclusion under an early-warning constraint: visible diffusion differences in fully observed cascades do not translate automatically into strong early veracity classification. The reach result generalizes more clearly than the veracity result.

This does not mean the earlier literature was wrong. Rather, it means that the inferential target is different. A retrospective cascade comparison asks whether false and true content differ once full trajectories are observed. The present thesis asks whether enough of those differences are already visible at the very beginning of diffusion to support prediction under a real early-warning constraint. The answer is more favorable for reach than for veracity. In that sense, the current results partly support the descriptive literature on veracity-linked diffusion differences, but they also qualify it by showing that those differences are much harder to recover from partial early observations than from completed cascades.

The results also fit reasonably well with broader work on early popularity and diffusion forecasting. Studies such as \citet{szabo_predicting_2010} and \citet{cheng_can_2014} show that early traces can be informative about later aggregate outcomes, especially when the target is eventual scale rather than semantic truth. My results are closer to that line of evidence on the reach side. What is more surprising is that the same early diffusion information does not transfer as strongly to veracity, even though veracity-linked diffusion differences are prominent in the misinformation literature. That is why the task asymmetry is not just a numerical pattern; it is the substantive contribution of the thesis.

\subsection{Why the Methodological Changes Matter}
The residualization diagnostics (Results Table~\ref{tab:residualization_diag}) indicate that structural features were meaningfully volume-correlated before adjustment and nearly orthogonal to log early volume after adjustment. This matters because, without adjustment, apparent structure effects can be volume effects in disguise. Given how strongly early volume predicts later reach, that distinction is important.

This is also where the thesis connects back to \citet{goel_structural_2016}. If structural quantities are mechanically tied to cascade size, then any predictive success can be hard to interpret. The residualization step does not solve that problem completely, but it does make the feature bundles more defensible by separating shape from early volume as much as the observed data allow.

Primary reliance on $W_k$ windows also remains important for interpretation. By fixing observed size, these analyses reduce confounding by raw arrival speed and make comparisons across cascades more structurally meaningful. The time-window results remain useful, but they answer a slightly different question. They show what can be recovered after a fixed amount of elapsed time, not after a fixed amount of observed information.

This point is especially important for how the thesis is positioned methodologically. The thesis does not claim that time windows are uninformative or illegitimate. Instead, it shows that the two window definitions answer slightly different empirical questions. A fixed-time design is closer to an operational clock, while a fixed-size design is closer to a controlled comparison of early observed information. Because the present thesis is mainly interested in what diffusion structure contributes beyond raw speed and volume, the $k$-window design is more defensible as the main specification. The weaker time-window results, especially for reach, strengthen rather than weaken that interpretation.

\subsection{Modeling Implications}
The model-family comparisons point to a similar task split. For veracity, logit remains competitive in the main FibVID analysis, and the gains from greater model flexibility are limited. This is consistent with the broader pattern in the thesis: the veracity signal is present, but weak enough that more flexible models do not transform the task.

For reach, the situation is different. Tree-based regressors become very strong in larger windows, which suggests that non-linear interactions matter more for diffusion-outcome prediction than for misinformation-risk screening. In practical terms, this implies two different modeling regimes. For early veracity screening, relatively simple probabilistic models may already capture most of the available signal. For reach forecasting, richer interactions and nonlinearities are more important.

The supplementary native-graph result adds one more modeling implication. For veracity, preserving more of the original claim-level topology appears to recover somewhat more signal than the harmonized main specification. I do not treat that result as a replacement for the main design, because the thesis still needs one common cascade-compatible framework across FibVID and Twitter15/16. But it does suggest that veracity may depend on more local or multi-origin topological patterns than reach does. Reach prediction benefits from a cleaner hierarchical cascade representation, whereas veracity appears more sensitive to how the original interaction structure is preserved.

\subsection{Strength of Evidence and Inference Boundaries}
The permutation tests make the task asymmetry sharper. On FibVID at $k=60$, observed reach $R^2$ is far above the shuffled-label null, while veracity AUC is only modestly separated from it. The Twitter15/16 benchmark shows the same qualitative pattern, although the veracity signal is somewhat stronger there. These checks make it difficult to read the veracity results as a strong classification success.

At the same time, the time-window results do not reverse the main pattern. They are weaker than the size-based results, especially for reach, but they still point in the same direction: reach is more predictable than veracity, and $k$-windows remain the stronger main specification. I therefore interpret the time-window analysis as a robustness check on window definition rather than as a competing primary design.

\subsection{Limitations and Threats to Validity}
\paragraph{Representation dependence.}
The main FibVID analysis depends on a harmonized cascade-compatible representation. This is a transparent and rule-based conversion, but it is still a representation choice rather than a claim that FibVID is natively identical to Twitter15/16. The main results should therefore be interpreted as evidence under that representation. The supplementary graph-native veracity check is useful here: it shows that some conclusions, especially on the veracity side, are sensitive to how propagation structure is represented.

\paragraph{Predictive, not causal.}
Even after residualization, diffusion features can proxy unobserved event-level shocks, platform effects, or coordinated behavior. These estimates are predictive associations, not causal effects of structure on misinformation or reach.

\paragraph{Finite-fold uncertainty.}
Inference uses fixed cross-validation folds. Some late-window reach estimates are very strong, and while this is substantively informative, their exact magnitude should not be overinterpreted. The thesis is more credible on the direction and relative strength of the effects than on any single point estimate.

\paragraph{Representation limits of interaction traces.}
Both the released rumor trees in Twitter15/16 and the propagation records in FibVID are interaction traces rather than full exposure networks. They do not observe recommendation systems, passive exposure, or off-platform diffusion. The models therefore recover signal from visible propagation, not from the full information environment in which these cascades spread.

\paragraph{Dataset composition.}
FibVID also differs from Twitter15/16 in ways beyond raw file format. It is a more claim-centered dataset, the main binary label setup is more imbalanced, and the harmonized representation is built from a different released structure than the native Twitter rumor trees. These differences do not invalidate comparison, but they do mean that cross-dataset contrasts should be read as evidence about portability under different data environments rather than as a pure apples-to-apples test of one fixed benchmark.

\subsection{Implications for Early Warning Systems}
The practical implication is a staged early-warning logic. For veracity, diffusion-only performance remains limited in both datasets, so the most defensible use is early screening or triage rather than confident truth classification. Such signals may still help prioritize which cascades deserve further verification, but they should be interpreted conservatively.

For reach, the case is much stronger. As more early interactions accumulate, structural and dynamic features become highly informative about eventual spread. This makes diffusion-based models more naturally suited to prioritizing likely high-impact cascades than to deciding whether a claim is true or false.

Methodologically, the main implication is that diffusion-based early warning should not be framed as a single undifferentiated problem. The results suggest that task, window definition, and structural representation all matter. In this thesis, the most defensible main design is a size-based, cascade-compatible specification for the reach task, with more cautious interpretation on the veracity side.

There is also a narrower practical implication for platform or monitoring settings. If a system has only a small amount of very early propagation information, the present evidence suggests it should use diffusion signals mainly to prioritize which cascades are likely to become large, not to make strong claims about truth status. Veracity-related diffusion features may still be useful as one input among others, but on their own they are not strong enough here to justify confident intervention. This distinction matters because the costs of a false positive in truth classification are different from the costs of prioritizing a cascade for further review.

Finally, the results suggest a modest but useful research agenda for diffusion-based misinformation work. One direction is to combine early diffusion with content or source-side features, since the present thesis shows that diffusion alone leaves veracity underidentified. Another is to test whether the reach/veracity asymmetry persists in newer or platform-specific propagation settings. A third is to develop more native representations for datasets like FibVID without giving up too much cross-dataset comparability. These are extensions of the present findings, not open holes in the current design.
